\BourbakiExercises{习题}

\begin{enumerate}
\def\labelenumi{\arabic{enumi})}
\tightlist
\item
  设 \(\mathrm{K}\) 是一个特征为 \(p\) 的交换域，\(\mathrm{M}\) 是一个 \(p\)-准素 \(\mathbf{Z}\)-模，\(\mathrm{A}\) 是群 \(\mathrm{M}\) 在 \(\mathrm{K}\) 上的代数，\(\mathfrak{m}\) 是元素 \(\sum_{m} a_{m} e_{m}\)（属于 \(\mathrm{A}\)，参见 III, §2, No.~6, 第 446 页）中满足 \(\sum_{m} a_{m} = 0\) 者所成的集。
\end{enumerate}

\begin{enumerate}
\def\labelenumi{\alph{enumi})}
\item
  证明 A 是一个局部环，其极大理想为 \(\mathfrak{m}\)，并且 \(\mathfrak{m}\) 是一个诣零理想（参见 VIII, 第 41 页，习题 10）。
\item
  此外，假设位似 \(p_{\mathrm{M}}\) 是满射。证明我们有 \(\mathfrak{m} = \mathfrak{m}^2\)。
\end{enumerate}

\begin{enumerate}
\def\labelenumi{\arabic{enumi})}
\setcounter{enumi}{1}
\item
  设 A 是一个环，M 是一个无根 A-模。证明位似环 \(A_{M}\) 无根（把 \(A_{s}\) 与 \(M^{M}\) 的一个子模等同起来）。
\item
  \begin{enumerate}
  \def\labelenumii{\alph{enumii})}
  \tightlist
  \item
    设 A 是一个环，其商环 \(\mathrm{A} / \Re (\mathrm{A})\) 是半单的。证明对每个 A-模 M，我们有 \(\Re (\mathrm{M}) = \Re (\mathrm{A})\mathrm{M}\)（注意 \(\mathrm{M} / \Re (\mathrm{A})\mathrm{M}\) 是一个 \(\mathrm{A} / \Re (\mathrm{A})\)-模）。
  \end{enumerate}
\end{enumerate}

\begin{enumerate}
\def\labelenumi{\alph{enumi})}
\setcounter{enumi}{1}
\tightlist
\item
  由此给出一个环 A 和一个 A-模的族 \((\mathrm{V}_{i})_{i\in\mathrm{I}}\) 的例子，使得 \(\Re\big(\prod\mathrm{M}_{i}\big)\neq\prod\Re(\mathrm{M}_{i})\)（对一个根不是有限生成的环利用 a)）。
\end{enumerate}

\begin{enumerate}
\def\labelenumi{\arabic{enumi})}
\setcounter{enumi}{3}
\item
  设 A 是一个交换域 K 上的代数，a 是 A 的一个双边理想。设 B 是由 1 与 a 生成的 A 的子代数。证明 B 的根含于 A 的根之中（注意若 S 是一个单 A-模，则或者对每个 \(x \in S\) 有 ax = 0，或者对 S 中每个 \(x \neq 0\) 有 ax = S）。
\item
  设 \(A\) 是一个交换域 \(K\) 上的代数。
\end{enumerate}

\begin{enumerate}
\def\labelenumi{\alph{enumi})}
\tightlist
\item
  证明 A 的根中在 K 上为代数元的元素是幂零的。b) 假设 K 是一个无限域且 \(\operatorname{Card}(K) > [A : K]\)。证明 A 的根是一个诣零理想（按 VIII, 第 47 页的引理 1 的方式论证）。
\end{enumerate}

\begin{enumerate}
\def\labelenumi{\arabic{enumi})}
\setcounter{enumi}{5}
\tightlist
\item
  设 \(\mathrm{A}\) 是一个环。
\end{enumerate}

\begin{enumerate}
\def\labelenumi{\alph{enumi})}
\item
  设 \(\mathfrak{a}\) 与 \(\mathfrak{b}\) 是 A 的双边理想。若 \(\mathfrak{a}\) 与 \(\mathfrak{b}\) 都是幂零的（相应地，诣零理想），则 \(\mathfrak{a} + \mathfrak{b}\) 亦然。
\item
  证明所有双边诣零理想之和 \(\mathfrak{G}\) 是一个双边诣零理想，而所有幂零双边理想之和 \(\mathfrak{N}\) 是一个包含所有幂零理想的双边诣零理想（VIII, 第 149 页，习题 3）。
\item
  证明 G 是 A 的最大半素诣零理想（同前所引）；我们称 G 为 A 的幂零根。若 A 是交换的，则 G 与 N 都等于 A 的幂零元素所成之集。
\end{enumerate}

¶ 7) 设 A 是一个环，I 是这样的序数 \(\alpha\) 所成之集：\(\text{Card}(\alpha) \leqslant \text{Card}(A)\)（《集合论》，III, §6, 第 241 页，习题 10）。用超限归纳定义双边理想 \(\mathfrak{N}_{\alpha}\)（对每个 \(\alpha \in I\)）如下。取 \(\mathfrak{N}_0 = \mathfrak{N}\)（习题 6, b)）。若 \(\alpha\) 有后继 \(\alpha + 1\)，则 \(\mathfrak{N}_{\alpha+1}\) 是使 \(\mathfrak{N}_{\alpha+1}/\mathfrak{N}_{\alpha}\) 为 \(A/\mathfrak{N}_{\alpha}\) 的诸幂零双边理想之和的双边理想。若 \(\alpha\) 是一个极限序数，则取 \(\mathfrak{N}_{\alpha}\) 为诸 \(N_{\beta}\)（\(\beta < \alpha\)）的并。设 \(\tau\) 是使 \(N_{\tau+1} = N_{\tau}\) 的最小序数；理想 \(P = N_{\tau}\) 称为 A 的素根。

\begin{enumerate}
\def\labelenumi{\alph{enumi})}
\item
  证明我们有 \(N \subset P \subset G\)，且 P 是 A 的最小半素诣零理想。\(^{(1)}\)
\item
  证明 \(\mathfrak{P}\) 是 A 的诸半素双边理想的交（参见 VIII, 第 149 页，习题 3）。
\item
  证明 P 是 A 的诸素双边理想的交（同前所引，c)）。d) 证明 P 是 A 中具有下述性质的元素 a 所成之集：A 中元素的每个序列 \((a_{n})_{n\geqslant0}\)，若它满足 \(a_{0}=a\) 且对每个 n 有 \(a_{n+1}\in a_{n}Aa_{n}\)，则它具有有限支集（方法同同前所引）。
\end{enumerate}

\begin{enumerate}
\def\labelenumi{\arabic{enumi})}
\setcounter{enumi}{7}
\tightlist
\item
  我们沿用习题 7 的记号。设 A 是一个左诺特环。
\end{enumerate}

\begin{enumerate}
\def\labelenumi{\alph{enumi})}
\item
  证明 \(\mathfrak{N}\) 是 A 的最大幂零双边理想，且它等于 A 的素根 \(\mathfrak{P}\)。
\item
  证明 A 的每个左或右诣零理想都是幂零的（对环 A/N 应用 VIII, 第 150 页的习题 5, a)）。于是我们有 \(\mathfrak{N} = \mathfrak{P} = \mathfrak{G}\)。
\end{enumerate}

\begin{enumerate}
\def\labelenumi{\arabic{enumi})}
\setcounter{enumi}{8}
\item
  设 \(\mathrm{K}\) 是一个交换域，\(\mathrm{B}\) 是环 \(\mathrm{K}[(\mathrm{X}_n)_{n\geqslant 1}]\)，\(\mathfrak{I}\) 是 \(\mathrm{B}\) 中由序列 \((\mathrm{X}_n^n)\) 生成的理想，\(\mathrm{A}\) 是环 \(\mathrm{B} / \mathfrak{I}\)。证明 \(\mathrm{A}\) 的根是一个诣零理想但不是幂零的。
\item
  设 A 是一个环。
\end{enumerate}

\begin{enumerate}
\def\labelenumi{\alph{enumi})}
\item
  设 Z 是 A 的中心。证明我们有 \(Z \cap \Re(A) \subset \Re(Z)\)（参见下文习题 15, b)，其中给出该包含为严格包含的例子）。
\item
  证明 \(\Re(\mathrm{A})\) 不含非零幂等元。
\item
  证明 \(\Re(\mathrm{A})\) 与 A 的左底座 s（VIII, 第 130 页，习题 8）的交是一个平方为零的双边理想，且它是 s 与其左零化子的交（利用 VIII, 第 130 页的习题 7 与 8, a)）。
\end{enumerate}

\(d\) ) 设 \(e\) 是 A 中的一个幂等元。证明环 \(e\mathrm{A}e\) 的根是 \(e\Re (\mathrm{A})e\)（注意对 \(a\in e\mathrm{A}e\)，我们有 \((1 - a)(1 - e) = (1 - e)(1 - a) = 1 - e\)）。

\begin{enumerate}
\def\labelenumi{\alph{enumi})}
\setcounter{enumi}{4}
\tightlist
\item
  证明环 \(\mathbf{M}_{n}(\mathrm{A})\) 的根是理想 \(\mathbf{M}_{n}(\Re(\mathrm{A}))\)（VIII, 第 114 页，例 2）。
\end{enumerate}

\begin{enumerate}
\def\labelenumi{\arabic{enumi})}
\setcounter{enumi}{10}
\tightlist
\item
  设 A 是一个交换环，B 是环 A{[}X{]}。
\end{enumerate}

\begin{enumerate}
\def\labelenumi{\alph{enumi})}
\item
  首先假设 A 没有非零的幂零元素。设 \(f = \sum_{i} a_{i} X^{i}\) 与 \(g = \sum_{i} b_{i} X^{i}\) 是 B 的两个元素。证明我们有 \(a_{p} b_{q} = 0\)，只要 \(p + q > \deg(fg)\)。b) 设 N 是 A 的幂零根（习题 6）。证明多项式 \(f \in B\) 在 B 中可逆当且仅当其常数项在 A 中可逆且其余系数都属于 N（对 f 在 \((\mathrm{A}/\mathfrak{N})[\mathrm{X}]\) 中的像应用 a)）。
\item
  由此推出 B 的根是理想 \(\Re[X]\)，且它等于 B 的幂零根。
\end{enumerate}

\begin{enumerate}
\def\labelenumi{\arabic{enumi})}
\setcounter{enumi}{11}
\tightlist
\item
  \begin{enumerate}
  \def\labelenumii{\alph{enumii})}
  \tightlist
  \item
    环 A 无根当且仅当它同构于本原环的积 \(\prod_{i\in I}B_{i}\) 的一个子环 C，使得 \(\mathrm{pr}_{i}(\mathrm{C})=\mathrm{B}_{i}\) 对每个 \(i\in I\) 成立（考虑单 A-模的类所成之集）。
  \end{enumerate}
\end{enumerate}

\begin{enumerate}
\def\labelenumi{\alph{enumi})}
\setcounter{enumi}{1}
\tightlist
\item
  设 V 是体 D 上的一个无限维向量空间，B 是环 \(\operatorname{End}_{\mathrm{D}}(\mathrm{V})\)，s 是它的底座（参见 VIII, 第 130 页，习题 8）。设 C 是 \(B \times B\) 的由 \(s \times s\) 与单位元生成的子环。证明 C 无根但不同构于本原环的一个积。
\end{enumerate}

\begin{enumerate}
\def\labelenumi{\arabic{enumi})}
\setcounter{enumi}{12}
\item
  设 \(G\) 是一个无挠交换群，\(K\) 是一个交换域。证明代数 \(K[G]\) 无根（给 \(G\) 装备一个与其群结构相容的全序，参见 VI, §1, 第 32 页，习题 20）。
\item
  设 \(\mathrm{K}\) 是一个序交换域，\(\mathrm{G}\) 是一个群。证明群 \(\mathrm{G}\) 在 \(\mathrm{K}\) 上的代数不具有非零的（左或右）诣零理想。（对 \(x = \sum a_{g}e_{g}\)，令 \(t(x) = a_{1}\)，\(x^{*} = \sum a_{g}e_{g - 1}\)。注意若 \(x\) 非零，则 \(t((xx^{*})^{2k}) > 0\) 对每个 \(k\geqslant 0\) 成立。）对群 \(\mathrm{G}\) 在体 \(\mathrm{K}(i)\)（其中 \(i^2 = -1\)）上的代数证明同样的性质。
\item
  设 \(V_{0}\) 是体 D 上的一个有限维 n 向量空间，B 是环 \(\operatorname{End}_{\mathrm{D}}(\mathrm{V}_{0})\) 的一个子环。令 \(V_{k}=V_{0}\)（\(k\geqslant1\)），并令 \(V=\oplus_{k\geqslant0}V_{k}\)。考虑 V 的满足下列条件的自同态 u 所成之集 A：存在一个（依赖于 u 的）整数 m，使得 \(u(V_{k})\subset\oplus_{i\leqslant m}V_{i}\) 对 \(k\leqslant m\) 成立，\(u(V_{k})\subset V_{k}\) 对 k\textgreater m 成立，且对 k\textgreater m，由限制得到的 \(V_{k}\) 的自同态与 k 无关并属于 B。
\end{enumerate}

\begin{enumerate}
\def\labelenumi{\alph{enumi})}
\tightlist
\item
  证明 A 是一个本原环，且其底座由 A 中使得对除有限多个指标 k 外都有 \(u(\mathrm{V}_{k}) = 0\) 的元素 u 组成（VIII, 第 130 页，习题 8）。
\item
  选取适当的 D、n 与 B，给出一个具有非零幂零根的商环的本原环的例子，以及一个中心有非零根的本原环的例子。
\end{enumerate}

\begin{enumerate}
\def\labelenumi{\arabic{enumi})}
\setcounter{enumi}{15}
\tightlist
\item
  设 A 是一个环，\(\mathfrak{a}\) 是含于 A 的根中的一个双边理想。设 M 与 N 是 A-模，\(u: \mathrm{M} \to \mathrm{N}\) 是一个同态。用 \(\overline{u}: \mathrm{M} / \mathfrak{aM} \to \mathrm{N} / \mathfrak{aN}\) 表示 A/ \(\mathfrak{a}\)-模的、由 \(u\) 导出的同态。
\end{enumerate}

\begin{enumerate}
\def\labelenumi{\alph{enumi})}
\item
  假设 M 与 N 都是有限生成的且 N 是投射的。若 \(\overline{u}\) 是双射，则 u 亦然。
\item
  给出一个 M 与 N 都是秩 1 自由模且 \(\overline{u}\) 是单射但 \(\operatorname{Ker}(u) \neq 0\) 的例子。c) 假设 M 与 N 都是投射的。若 \(\overline{u}\) 诱导出到一个直因子上的同构，则 u 是单射（化归到 M 与 N 是自由且相等并且 \(\overline{u}\) 是恒等的情形，然后利用 a)）。此外，若 M 是有限生成的，则 u 的像是 N 的一个直因子。
\item
  假设 M 是有限生成的且 N 是投射的。若 \(\overline{u}\) 是双射，则 u 亦然（把 c) 应用于一个同态 \(v: N \to M\)，它满足 \(\overline{v} = \overline{u}^{-1}\)，从而断定 N 是有限生成的）。
\item
  若 M 是投射的且 M/ \(\mathfrak{aM}\) 为零，则 M 为零。
\end{enumerate}

\begin{enumerate}
\def\labelenumi{\arabic{enumi})}
\setcounter{enumi}{16}
\tightlist
\item
  设 \(\mathbf{K}\) 是一个交换域，\(\mathbf{L}\) 是一个集。
\end{enumerate}

\begin{enumerate}
\def\labelenumi{\alph{enumi})}
\tightlist
\item
  证明存在一个 K-代数 A，它有一个由单位元、一个族 \((c_{\lambda})_{\lambda \in \mathrm{L}}\) 和一个族 \((r_{\lambda \mu})_{(\lambda ,\mu)\in \mathrm{L}\times \mathrm{L}}\) 组成的基，其乘法表为
\end{enumerate}

\[
\begin{array}{r} c _ {\lambda} ^ {2} = c _ {\lambda}, \qquad c _ {\lambda} c _ {\mu} = r _ {\lambda \mu} \quad \mathrm{if} \lambda \neq \mu \\ c _ {\lambda} r _ {\mu \nu} = \delta_ {\lambda \mu} r _ {\mu \nu}, \qquad r _ {\lambda \mu} c _ {\nu} = \delta_ {\mu \nu} r _ {\lambda \mu} \\ r _ {\lambda \mu} r _ {\nu \rho} = 0, \end{array}
\]

其中 \(\delta_{\mu\nu}\) 是克罗内克 δ 函数。证明 \(\Re(\mathrm{A})\) 是由诸 \(r_{\lambda\mu}\) 生成的 A 的子代数；我们有 \(\Re(\mathrm{A})^{2}=0\)。

\begin{enumerate}
\def\labelenumi{\alph{enumi})}
\setcounter{enumi}{1}
\tightlist
\item
  假设 \(\operatorname{Card}(\mathrm{L}) \geqslant 2\)。证明 A 的中心 Z 是由 1、诸 \(r_{\lambda\lambda}\)（\(\lambda \in L\)）以及（若 L 有限）元素 \(\sum_{\lambda} c_{\lambda} - \sum_{\mu \neq \nu} r_{\mu\nu}\) 生成的 K-向量空间。由此推出，对 \(\lambda \in L\)，Z 中不存在与 \(c_{\lambda} \pmod{\Re(A)}\) 同余的幂等元。
\end{enumerate}

¶ 18) a) 设 A 是一个环，a 是 A 的一个双边诣零理想，\((\overline{e}_{n})_{n\geqslant1}\) 是 A/a 中的一个正交幂等元序列。证明存在 A 中的一个正交幂等元序列 \((e_{n})_{n\geqslant1}\)，使得对每个 n，\(\overline{e}_{n}\) 是 \(e_{n}\) 在 A/a 中的类。（设 \(\overline{e}\) 与 \(\overline{e}'\) 是 A/a 中两个正交幂等元，\(e'\) 是 A 中提升 \(\overline{e}'\) 的幂等元，a 是 e 在 A 中的一个代表。注意我们可以修改 a 使 \(ae' = e'a\)，然后如同 VIII, 第 159 页的命题 7 那样使 \(a^{2} = a\)。）

\begin{enumerate}
\def\labelenumi{\alph{enumi})}
\setcounter{enumi}{1}
\tightlist
\item
  设 K 是一个交换域，L 是一个不可数集，A 是上一习题中定义的 K-代数。证明 A 中不存在正交幂等元族 \((e_{\lambda})_{\lambda\in\mathrm{L}}\)，使得 \(e_{\lambda}\equiv c_{\lambda} \pmod{\Re(\mathrm{A})}\) 对每个 \(\lambda\in L\) 成立。（考虑由元素 \((\lambda,\mu)\)（属于 \(L\times L\)）中满足 \(c_{\lambda}(e_{\mu}-c_{\mu})=0\) 者组成的集 H。注意诸集 \(H\cap(\{\lambda\}\times L)\) 对每个 \(\lambda\in L\) 都是有限的，而诸集 \(\mathbb{C}H\cap(L\times\{\mu\})\) 对每个 \(\mu\in L\) 都是有限的。）
\end{enumerate}

\begin{enumerate}
\def\labelenumi{\arabic{enumi})}
\setcounter{enumi}{18}
\item
  设 A 是一个环，a 是含于 \(\Re(\mathrm{A})\) 中的一个双边理想，e 与 \(e'\) 是 A 中的幂等元，\(\overline{e}\) 与 \(\overline{e}'\) 是它们在环 \(\overline{A} = A/a\) 中的类。证明若幂等元 \(\overline{e}\) 与 \(\overline{e}'\) 等价（VIII, 第 39 页，习题 5），则 e 与 \(e'\) 亦然（注意 Ae 是 A-模 \(\overline{A}\overline{e}\) 的一个投射覆盖）。
\item
  设 \(\mathrm{A}\) 是交换域 \(\mathrm{K}\) 上的一个代数，使得 \(\Re(\mathrm{A})\) 是一个诣零理想，且代数 \(\mathrm{A} / \Re(\mathrm{A})\) 同构于 \(\mathrm{K}\) 上矩阵代数的一个有限积。证明存在一个子代数 \(\mathrm{S}\)（在 \(\mathrm{A}\) 内），使 \(\mathrm{K}\)-向量空间 \(\mathrm{A}\) 是 \(\mathrm{S}\) 与 \(\Re(\mathrm{A})\) 的直和（利用上面的习题 18 与 19，以及 VIII, 第 94 页的习题 19）。
\item
  设 A 是一个环，\(u: P \to M\) 是 A-模的一个满同态，N 是它的核。用 \(\operatorname{End}_{\mathrm{A}}(\mathrm{P}; \mathrm{N})\)（相应地，\(\operatorname{Aut}_{\mathrm{A}}(\mathrm{P}; \mathrm{N})\)）表示 P 的满足 \(u(\mathrm{N}) \subset \mathrm{N}\) 的自同态 u 所成之环（相应地，自同构 u 所成之群）。
\end{enumerate}

\begin{enumerate}
\def\labelenumi{\alph{enumi})}
\tightlist
\item
  定义一个正合列
\end{enumerate}

\[
0 \to \mathrm{Hom} _ {\mathrm{A}} (\mathrm{P}, \mathrm{N}) \xrightarrow {i} \mathrm{End} _ {\mathrm{A}} (\mathrm{P}; \mathrm{N}) \xrightarrow {p} \mathrm{End} _ {\mathrm{A}} (\mathrm{M})
\]

其中 \(i\) 的像是一个双边理想 \(\mathfrak{I}\)（环 \(\operatorname{End}_{\mathrm{A}}(\mathrm{P};\mathrm{N})\) 的）。

\begin{enumerate}
\def\labelenumi{\alph{enumi})}
\setcounter{enumi}{1}
\tightlist
\item
  假设 \((P, u)\) 是 M 的一个投射覆盖。则 p 是满射，且诱导出一个满同态 \(p' : \operatorname{Aut}_{\mathrm{A}}(\mathrm{P}; \mathrm{N}) \to \operatorname{Aut}_{\mathrm{A}}(\mathrm{M})\)。我们有群的一个正合列
\end{enumerate}

\[
0 \to 1 _ {\mathrm{P}} + \Im \to \mathrm{Aut} _ {\mathrm{A}} (\mathrm{P}; \mathrm{N}) \xrightarrow {p ^ {\prime}} \mathrm{Aut} _ {\mathrm{A}} (\mathrm{M}) \to 0;
\]

若 P 是自由的且 N 不约化为 0，则同态 \(p'\) 不是单射。

\begin{enumerate}
\def\labelenumi{\arabic{enumi})}
\setcounter{enumi}{21}
\tightlist
\item
  设 \(A\) 是一个交换整环，它只有有限多个 \(s > 1\) 个极大理想 \(\mathfrak{m}_1, \ldots, \mathfrak{m}_s\)。
\end{enumerate}

\begin{enumerate}
\def\labelenumi{\alph{enumi})}
\item
  证明环 A/ \(\Re(A)\) 同构于诸体 A/mi 的积。
\item
  证明 A-模 \(A/m_{i}\) 不具有投射覆盖，尽管它们作为 \(A/\Re(A)\)-模是投射的且有限生成的（注意 A 是 \(A/\Re(A)\) 的一个投射覆盖）。
\item
  设 K 是一个交换域。证明由 \(\mathrm{K(T)}\) 中可以写成 \(\mathrm{P(T)/Q(T)}\)（其中 \(\mathrm{Q(0) \neq 0}\) 且 \(\mathrm{Q(1) \neq 0}\)）的有理分式组成的子环 A 具有所需的性质。
\end{enumerate}

\begin{enumerate}
\def\labelenumi{\arabic{enumi})}
\setcounter{enumi}{22}
\item
  设 A 是一个环，a 是 A 的一个（左）理想。假设对 a 的元素的每个序列 \((a_{n})_{n\geqslant1}\)，存在一个整数 m 使得 \(a_{1}\cdots a_{m}=0\)。证明对每个非零 A-模 M，我们有 \(aM\neq M\)（否则，用归纳法构造 a 中的序列 \((a_{n})_{n\geqslant1}\) 与 M 中的序列 \((x_{n})_{n\geqslant1}\)，使得对每个 p 都有 \(a_{1}\cdots a_{p}x_{p}\neq0\)）。
\item
  设 \(A\) 是一个环，\(\mathfrak{r}\) 是它的根。证明下列性质等价：
\end{enumerate}

\begin{enumerate}
\def\labelenumi{(\roman{enumi})}
\tightlist
\item
  每个有限生成左 A-模都具有一个投射覆盖。
\end{enumerate}

(i') 每个有限生成右 A-模都具有一个投射覆盖。

\begin{enumerate}
\def\labelenumi{(\roman{enumi})}
\setcounter{enumi}{1}
\tightlist
\item
  环 \(\mathrm{A} / \mathfrak{r}\) 是半单的，且 \(\mathrm{A} / \mathfrak{r}\) 中的每个幂等元都是 \(\mathrm{A}\) 中某个幂等元的类。
\end{enumerate}

（为证 (i)⇒(ii)，注意理想的每个直和分解 \(A/r = a_{1} \oplus a_{2}\) 都可提升为分解 \(A = P_{1} \oplus P_{2}\)，其中 \(P_{i}\) 是 \(a_{i}\) 的一个投射覆盖。此外，证明 (ii) 蕴含每个 A/r-模都具有形式 \((\mathrm{A}/\mathfrak{r}) \otimes_{\mathrm{A}} \mathrm{P}\)，其中 P 是一个投射 A-模。）

¶ 25) 设 A 是一个环，r 是它的根。证明下列条件等价：

\begin{enumerate}
\def\labelenumi{(\roman{enumi})}
\tightlist
\item
  每个左 A-模都具有一个投射覆盖。
\end{enumerate}

(i') 每个右 A-模都具有一个投射覆盖。

\begin{enumerate}
\def\labelenumi{(\roman{enumi})}
\setcounter{enumi}{1}
\tightlist
\item
  环 A/r 是半单的，且对 r 的元素的每个序列 \((a_{n})_{n\geqslant1}\)，存在一个整数 m 使得 \(a_{1}\cdots a_{m}=0\)。
\end{enumerate}

（用习题 24 中相应蕴含关系的同样方式证明 (ii)⇒(i)，并利用习题 23。假设条件 (i) 成立。考虑 A \(^{(N)}\) 的由 \(e_1 - a_1e_2\)、\(e_2 - a_2e_3\)、\ldots{} 生成的子模 M。把习题 16, e) 应用于 A \(^{(N)}\) /M 的一个投射覆盖，推出 M = A \(^{(N)}\)，这蕴含当 n 充分大时 \(a_1 \cdots a_n = 0\)。）

\begin{enumerate}
\def\labelenumi{\arabic{enumi})}
\setcounter{enumi}{25}
\tightlist
\item
  若环 A 的每个非诣零左理想都含有一个非零幂等元，则称 A 为一个佐恩环。
\end{enumerate}

\begin{enumerate}
\def\labelenumi{\alph{enumi})}
\item
  证明在佐恩环中，每个非诣零右理想都含有一个非零幂等元（注意对 A 的每个非幂零元素 \(a\)，存在一个 \(x \in A\) 使得 \(xaxa = xa \neq 0\)）。
\item
  证明佐恩环的根是一个诣零理想，且无零因子的佐恩环是一个体。
\item
  证明底座非零的本原环是一个佐恩环（参见 VIII, 第 129 页，习题 4, b) 及第 130 页，习题 7, a)）。给出一个不是佐恩环的本原环的例子（参见 VIII, 第 132 页，习题 13），以及一个底座非零但中心不是佐恩环的本原环的例子（利用习题 12 的方法）。
\end{enumerate}

\(d\) ) 设 K 是一个交换域，A 是 \(\mathrm{K}(\mathrm{X})^{\mathbf{N}}\) 中由满足下列条件的序列 \(s = (f_n)\) 组成的子环：存在一个多项式 \(\varphi(s) \in \mathrm{K}[\mathrm{X}]\)，使得 \(f_n\) 等于 \(\varphi(s)\)（对充分大的 \(n\)）。证明 A 是一个交换佐恩环，无根，含有既不可逆又不是零因子的元素，但其像 \(\varphi(A)\) 不是佐恩环。

¶ 27) 若对 A 的每个元素 a，存在一个 \(x \in A\) 使得 axa = a，则称环 A 是绝对平坦的，或冯·诺伊曼正则的。

\begin{enumerate}
\def\labelenumi{\alph{enumi})}
\item
  环 A 是绝对平坦的当且仅当 A 的每个单生成左（或右）理想在 \(A_{s}\) 中有一个补（也就是说，由一个幂等元生成）。
\item
  证明绝对平坦环的每个有限生成理想都由一个幂等元生成。（化归到理想 \(Ae + Af\) 的情形，其中 e 与 f 是幂等元。然后，把 Ae 换成 \(Ae(1 - f)\)，化归到 ef = 0 的情形，并考虑元素 \(e + f - fe\)。）（左或右）诺特绝对平坦环是半单的。
\item
  由 b) 推出，绝对平坦环的两个单生成左（相应地，右）理想的交是单生成的（考虑右零化子）。
\item
  绝对平坦环的每个商环都是绝对平坦的，绝对平坦环的每个积亦然。
\item
  证明绝对平坦环的根约化为 0。由此推出 A 的每个双边理想都是包含它的诸极大左理想的交。
\end{enumerate}

\(f\) ) 绝对平坦环的中心是绝对平坦的（证明若 \(a \in \mathbb{Z}\) 与 \(x \in \mathrm{A}\) 满足 \(a^2 x = a\)，则 \(a^2 x^n \in \mathbb{Z}\) 对每个 \(n \geqslant 0\) 成立，且 \(a(a^2 x^3)a = a\)）。\(g\) ) 证明向量空间的自同态环是绝对平坦环。

¶ 28) 设 A 是一个绝对平坦环（习题 27），n 是一个整数。

\begin{enumerate}
\def\labelenumi{\alph{enumi})}
\item
  证明 \(A^{n}\) 的每个有限生成子模都具有一个补子模。（对 n 用归纳法论证，先证明 \(M \cap A^{n-1}\) 是有限生成的，然后选取 \(M \cap A^{n-1}\) 在 \(A^{n-1}\) 中的补子模以及 M 的像在 \(Ae_{n}\) 中的补子模。）
\item
  由此推出环 \(\mathbf{M}_n(\mathrm{A})\) 是绝对平坦的（参见 VIII, 第 114 页，例 2）。
\end{enumerate}

\begin{enumerate}
\def\labelenumi{\arabic{enumi})}
\setcounter{enumi}{28}
\item
  设 V 是交换域 K 上的一个可数维向量空间，\(A = \text{End}_{\text{K}}(\text{V})\)。证明存在唯一的极大双边理想，即有限秩自同态所成之集 \(\text{End}_{\text{K}}^{\text{f}}(\text{V})\)（参见 VIII, 第 128 页的习题 2），但 A 的根为 0。
\item
  \begin{enumerate}
  \def\labelenumii{\alph{enumii})}
  \tightlist
  \item
    设 A 是一个除 0 外不含幂零元素的环。证明 A 中的每个幂等元 e 都属于 A 的中心（考虑元素 \((1-e)xe\) 与 \(ex(1-e)\)）。
  \end{enumerate}
\end{enumerate}

\begin{enumerate}
\def\labelenumi{\alph{enumi})}
\setcounter{enumi}{1}
\item
  设 A 是一个除 0 外不含幂零元素的绝对平坦环，a 是 A 的一个非零元素。证明存在 A 的一个元素 x，使得 e = ax = xa 是幂等的且 ea = ae = a。由此推出，对每个 \(y \in A\)，存在一个 \(z \in A\) 使得 ya = az，因而 A 的每个左或右理想都是双边的。
\item
  证明 A 的任何商环都不含非零幂零元素（利用 b)）。由此推出 A 同构于体的积 \(\prod_{\iota\in I}D_{\iota}\) 的一个子环 C，使得 \(\mathrm{pr}_{\iota}(\mathrm{C})=\mathrm{D}_{\iota}\) 对每个 \(\iota\in I\) 成立（由 b) 及 VIII, 第 129 页的习题 4 推出，无幂零元素的绝对平坦本原环是一个体，并利用上面的习题 12）。
\end{enumerate}

\(d\) ) 设 \((\mathrm{D}_{\iota})_{\iota \in \mathbb{I}}\) 是一个无限体族，B 是环 \(\prod_{\iota \in \mathrm{I}}\mathrm{D}_{\iota}\)。对 I 的任一子集 H，用 \(\mathrm{B_H}\) 表示 B 中由族 \((b_{\iota})_{\iota \in \mathrm{I}}\)（其中 \(b_{\iota} = 0\) 对 \(\iota \notin \mathrm{H}\)）组成的双边理想。设 \(\mathfrak{F}\) 是 I 的子集组成的一个递增有向集，它构成 I 的一个覆盖，A 是 B 的由 1 与诸 \(\mathrm{B_H}\)（\(\mathrm{H}\in \mathfrak{F}\)）生成的子环。证明 A 是绝对平坦的且不含非零幂零元素。

\begin{enumerate}
\def\labelenumi{\arabic{enumi})}
\setcounter{enumi}{30}
\tightlist
\item
  设 K 是一个交换域，A 是形如 \(\begin{pmatrix} a & b \\ 0 & c \end{pmatrix}\)（其中 \(a, b, c \in K\)）的矩阵所成之集。验证 A 是矩阵环的一个子环，且其根由 a = c = 0 确定。证明 \(\begin{pmatrix} 1 & 0 \\ 0 & 0 \end{pmatrix}\) 在 A/Re(A) 中的像是 A/Re(A) 中心的一个幂等元。
\end{enumerate}

